\documentclass[a4paper,12pt]{article}
\usepackage{fontspec}
\usepackage{amsmath,amssymb}
\usepackage{graphicx}
\usepackage[hidelinks]{hyperref}
\usepackage[margin=2.5cm]{geometry}
\usepackage[numbers]{natbib}
\usepackage{booktabs,array,float}
\title{From Compute to Complexity}
\author{Qinrang Liu}
\date{June 2026}
\begin{document}
\maketitle
\begin{abstract}
Test abstract.
\end{abstract}
\section{Introduction}
Some text.
\section{Results}
\subsection{The CST theorem}
We formalize the CST theorem on five axioms.
From these axioms:
\begin{equation}
CST = (S_c \cdot T_c) \cdot \exp(\alpha \cdot \Gamma_{st}) \tag{1}
\end{equation}
Spatial complexity Sc quantifies structural integration potential as the geometric mean of four orthogonal, MECE graph-theoretic measures:
\begin{equation}
S_c = (X_1 \cdot X_2 \cdot X_3 \cdot X_4)^{1/4} \tag{2}
\end{equation}
where C = global connectivity (LCC fraction); H = hierarchical depth.
\end{document}